Käesoleva lõputöö eesmärk oli uurida eDelivery dünaamilise otsingu mehhanismi 
ning luua selle põhjal töötav näidislahendus.

Loodud lahendus pidi võimaldama eFTI platvormi ja värava vahelist sõnumivahetust,
kasutades eDelivery SML- ja SMP-komponente saaja metaandmete dünaamiliseks leidmiseks.

Eesmärgi saavutamiseks analüüsiti esmalt olemasolevaid avatud lähtekoodiga 
AP- , SMP- ja SML-komponente ning hinnati nende sobivust.
Selle baasil arendati neljas etapis näidislahendus, liikudes lihtsaimast töötavast versioonist
täielikult funktsioneeriva dünaamilise sõnumivahetuseni.
Lõpetuseks teostati nii loodud näidislahendusele kui ka valmiskomponentidele lokaalses
testkeskkonnas jõudluspõhine hindamine ja tulemuste võrdlus.

Töö peamise tulemusena valmis toimiv sõnumivahetuse lahendus,
kus saatja juurdepääsupunkt suudab dünaamilise otsingu abil leida saaja metaandmed
ja edastada sõnumi teisele juurdepääsupunktile. Sellega täideti töö peamine tehniline eesmärk
ja tõestati eFTI regulatsioonis (2024/1942) kirjeldatud platvormi ja värava vahelise
dünaamilise suhtluse teostatavust \cite{eu_reg_2024_1942}.

Jõudlustestid kinnitasid, et loodud minimaalne SMP/SML implementatsioon suudab
väikese latentsusega töödelda suurt hulka päringuid. 

Näidislahendus on avaldatud GitHubi keskkonnas MIT litsentsi alusel \cite{github_edelivery_smp_sml}.

Töö oluliseks lisaväärtuseks on praktilise dokumentatsiooni loomine,
mis aitab leevendada eDelivery rakendusnäidete puudust ja vähendab tehnoloogia kasutuselevõtu
keerukust tulevikus.